\begin{abstract}
Parkinson's disease (PD) exhibits substantial heterogeneity in progression trajectories, yet current clinical stratification relies primarily on cross-sectional assessments. We applied unsupervised trajectory clustering to longitudinal motor and cognitive data from 536 early-stage PD patients followed over 5 years in the Parkinson's Progression Markers Initiative. Using variational autoencoder-based trajectory modeling and latent time alignment, we identified three robust progression subtypes: rapid motor decline (22\%, UPDRS-III slope 7.8$\pm$1.2 points/year), moderate progression (54\%, slope 2.9$\pm$0.6 points/year), and stable motor/cognitive decline (24\%, motor slope 1.2$\pm$0.6, MoCA slope -1.8$\pm$0.4 points/year). Subtypes exhibited distinct baseline predictors (baseline UPDRS, sex, age; classification accuracy 68\%, AUC 0.74) and differed significantly in biomarker profiles (DaTscan striatal binding, CSF $\alpha$-synuclein, GBA mutation prevalence). Simulated clinical trials demonstrated that enriching for rapid progressors reduced required sample size by 38\% while maintaining 80\% power, with effect size detectability improving from 2.8 to 1.7 UPDRS points/year. These data-driven subtypes provide a framework for precision prognostication and adaptive trial designs in PD.

\noindent\textbf{Keywords:} Parkinson's disease, Disease progression, Trajectory clustering, Clinical trials, Precision medicine
\end{abstract}
